	\begin{UseCase}{CU8}{Gestionar reglas de negocio}{
	Este caso de uso permite al analista visualizar los registros de las reglas de negocio del sistema.
	También permite al actor acceder a las operaciones de registro, consulta, modificación y eliminación de una regla de negocio.
	}
	\UCitem{Actor}{\hyperlink{jefe}{Líder de análisis}, \hyperlink{analista}{Analista}}
	\UCitem{Propósito}{Proporcionar al actor un mecanismo para llevar el control de las reglas de negocio de un proyecto.}
	\UCitem{Entradas}{Ninguna}
	\UCitem{Salidas}{\begin{itemize}
			\item \cdtRef{proyectoEntidad:claveProyecto}{Clave del proyecto}: Lo obtiene el sistema.
			\item \cdtRef{proyectoEntidad:nombreProyecto}{Nombre del proyecto}: Lo obtiene el sistema.
			\item \cdtRef{BREntidad}{Regla de negocio}: Tabla que muestra \cdtRef{BREntidad:claveBR}{clave}, y el \cdtRef{BREntidad:nombreBR}{Nombre} de todos los registros de las reglas de negocio.
			\item \cdtIdRef{MSG2}{No existe información}: Se muestra en la pantalla \IUref{IU9}{Gestionar Reglas de Negocio} cuando no existen reglas de negocio registradas.
	\end{itemize}}
	
	\UCitem{Precondiciones}{Ninguna}
	\UCitem{Postcondiciones}{Ninguna}
	\UCitem{Errores}{Ninguno}
	\UCitem{Tipo}{Primario}
\end{UseCase}
%--------------------------------------
\begin{UCtrayectoria}
	\UCpaso[\UCactor] Solicita gestionar las reglas de negocio seleccionando la opción ''Reglas de negocio'' del menú \IUref{MN3}{Menú de Proyecto}.
	\UCpaso[\UCsist] Obtiene la información de las reglas de negocio del proyecto seleccionado. \Trayref{GBR-A}
	\UCpaso[\UCsist] Muestra la información de las reglas de negocio en la pantalla \IUref{IU9}{Gestionar Reglas de negocio} y las operaciones disponibles de acuerdo a la regla de negocio \BRref{RN15}{Operaciones disponibles}.\label{CU8-P4}
	\UCpaso[\UCactor] Gestiona las reglas de negocio a través de los botones: \IUbutton{Registrar}, \editar , \eliminar y \raisebox{-1mm}{\includegraphics[height=11pt]{images/Iconos/consultar}}. 
\end{UCtrayectoria}		
%--------------------------------------
\hypertarget{CU8:TAA}{\textbf{Trayectoria alternativa A}}\\
\noindent \textbf{Condición:} No existen registros de Reglas de Negocio.
\begin{enumerate}
	\UCpaso[\UCsist] Muestra el mensaje \cdtIdRef{MSG2}{No existe información} en la pantalla \IUref{IU9}{Gestionar Reglas de negocio} para indicar que no hay registros de reglas de negocio para mostrar. \label{CU8-TA1}
	\UCpaso[\UCactor] Gestiona las reglas de negocio a través del botón: \IUbutton{Registrar}. 
	\item[- -] - - {\em {Fin del caso de uso}}.%
\end{enumerate}
%--------------

%--------------------------------------

\subsubsection{Puntos de extensión}

\UCExtenssionPoint{El actor requiere registrar una regla de negocio.}{Paso \ref{CU8-P4} de la trayectoria principal o del paso \ref{CU8-TA1} de la trayectoria alternativa A.}{\UCref{CU8.1}{Registrar Regla de negocio}}
\UCExtenssionPoint{El actor requiere modificar una regla de negocio.}{Paso \ref{CU8-P4} de la trayectoria principal.}{\UCref{CU8.2}{Modificar Regla de negocio}}
\UCExtenssionPoint{El actor requiere eliminar una regla de negocio.}{Paso \ref{CU8-P4} de la trayectoria principal.}{\UCref{CU8.3}{Eliminar Regla de negocio}}
\UCExtenssionPoint{El actor requiere consultar una regla de negocio.}{Paso \ref{CU8-P4} de la trayectoria principal.}{\UCref{CU8.4}{Consultar Regla de negocio}}